\chapter{Introduction}
\label{ch:introduction}

This report describes QLens, a locally hosted web application that takes an
untidy spreadsheet and turns it into a verified relational database in third
normal form (3NF). The application reads the data itself, works out which
columns determine which others, proposes a decomposition, and creates the
tables, keys and foreign keys in PostgreSQL. It also provides a place to query
the result and to study how those queries execute. This chapter explains the
problem that motivated the work, states the objectives, fixes the boundaries
of what was attempted, and walks through four concrete scenarios in which the
system is used.

%==============================================================
\section{Motivation and Problem Statement}
\label{sec:motivation}

\subsection{The fee register}

Consider a coaching institute with roughly three thousand fee payments
recorded over a year. The office keeps them in a single spreadsheet. Each row
holds the student's name and phone number, the batch the student attends, the
teacher of that batch together with the teacher's phone number, the month, the
amount paid and the mode of payment. The sheet is convenient because every fact
about a payment is visible on one line, and nobody had to design anything to
start using it.

The trouble appears after a few months. The same sheet now stores each
teacher's phone number once for every payment of every student in that
teacher's batches. When a teacher changes number, the clerk must find and edit
hundreds of rows, and any row that is missed leaves the sheet contradicting
itself. A student whose name was typed as \emph{Riya Sharma} in one month and
\emph{riya  sharma} in another is counted as two people, so the totals for
that student are split. Finally, a question as ordinary as ``who has not paid
for October?'' cannot be answered by any single filter, because the answer is
the set of students who are \emph{absent} from the October rows. In practice
the clerk works it out by hand, which takes an hour and is easy to get wrong.

None of these problems is a software defect in the spreadsheet program. They
are structural consequences of storing several different kinds of facts in one
table, and relational database theory has had precise names for them for
fifty years.

\subsection{The anomalies in the language of normalization}

Codd introduced the relational model and the idea of removing redundancy by
decomposing relations into normal forms~\cite{codd1970}. The standard
treatments~\cite{silberschatz2019,dbbook} classify the damage done by
redundancy into three anomalies, all of which the fee register exhibits.

\begin{itemize}
  \item \textbf{Update anomaly.} A single fact is stored in many rows, so a
    change must be applied to all of them. Teacher phone numbers repeat in
    every payment row of every student of that teacher.
  \item \textbf{Insertion anomaly.} A fact cannot be recorded until an
    unrelated fact exists. A new teacher with no enrolled student, or a new
    batch with no payment yet, has nowhere to live in the sheet.
  \item \textbf{Deletion anomaly.} Removing one fact destroys another. Deleting
    the only payment row of the last student in a batch also deletes the
    record that the batch exists and who teaches it.
\end{itemize}

These anomalies are explained by \emph{functional dependencies} (FDs). In the
register, a student determines a phone number, a student determines a batch, a
batch determines a teacher, and a teacher determines a phone number. Writing
this with the notation used throughout the report, the four dependencies are
\fd{student\_name}{student\_phone}, \fd{student\_name}{batch},
\fd{batch}{teacher} and \fd{teacher}{teacher\_phone}. The key of the sheet as a
whole is the pair (student, month), since a student pays once per month. Both
\code{teacher} and \code{teacher\_phone} therefore depend on only \emph{part}
of that key, and \code{teacher\_phone} depends on it only \emph{transitively},
through the batch and the teacher. A relation in which a non-key attribute
depends on part of a key violates second normal form, and one with a
transitive dependency violates third normal form. Decomposing the relation so
that every determinant is a key of its own table removes the anomalies.

Textbook exercises hand the student the FDs and ask for the decomposition. A
real spreadsheet comes with no such list. The dependencies have to be
discovered from the data, and the data is dirty. In the sample register that
accompanies this project, eight rows carry a spelling variant of a student's
name, and one teacher's old phone number survives in five rows after the
number was changed. A strict test would reject
\fd{student\_name}{student\_phone} because of those rows, even though the
intended dependency clearly holds. A usable tool must therefore distinguish
a dependency that is absent from one that is \emph{almost} present, and report
the offending rows to the user as inconsistencies to be repaired.

\subsection{Why existing tools do not close the gap}

Spreadsheet-like database products such as Airtable, NocoDB and Baserow give
a database look and feel but leave the schema design entirely to the user, who
is exactly the person that does not know what a functional dependency is.
General-purpose chat tools can answer questions about an uploaded file, but
they leave no lasting, constrained, correct database behind, and their answers
are not backed by keys or foreign keys. Academic normalization calculators
accept FDs typed by hand and never touch data. No widely available tool takes
a raw sheet and returns a normalized schema with the data already moved and
checked. QLens is an attempt to supply that missing step, with the database
engine itself doing as much of the analysis as possible.

%==============================================================
\section{Objectives}
\label{sec:objectives}

The overall aim is to build a working system that converts a flat spreadsheet
into a verified 3NF PostgreSQL database and makes the result queryable and
explainable. The aim is broken into the following objectives.

\begin{enumerate}
  \item Import a CSV or XLSX file into a staging table and infer an SQL type,
    a null count and a distinct count for every column.
  \item Discover functional dependencies and candidate keys from the data
    \emph{inside the database}, using grouped aggregation in SQL, and
    tolerate a small proportion of violating rows so that dependencies broken
    only by data-entry errors are still found.
  \item Detect near-duplicate text values using trigram similarity in
    PostgreSQL, and let the user apply fixes with a recorded audit trail.
  \item Compute a 3NF decomposition with a lossless-join guarantee, by
    synthesis from a minimal cover~\cite{bernstein1976}, and present it with
    plain-language reasons and generated DDL.
  \item Apply the accepted design in a single transaction, creating tables,
    primary keys, foreign keys and uniqueness constraints, migrating the data,
    and verifying that joining the new tables reproduces the original rows.
  \item Provide a sandboxed, read-only query workspace with saved queries and a
    complete run history.
  \item Provide a performance laboratory that exposes the real
    \code{EXPLAIN ANALYZE} plan and index suggestions, and a visualizer that
    steps through a \code{SELECT} stage by stage.
  \item Keep user data separate: every account owns workspaces, and every
    workspace is confined to its own schemas under its own database role.
\end{enumerate}

A second, equally important objective concerns the report itself. The design
deliberately relies on database features rather than application code
wherever possible: \code{CHECK}, \code{UNIQUE} and foreign key constraints,
generated columns, partial unique indexes, triggers, views with window
functions, PL/pgSQL functions, GIN indexes, and role-based privileges. The
metadata schema of the application is itself a normalized design, and the
later design chapters analyse it.

%==============================================================
\section{Scope of the Work}
\label{sec:scope}

The system is intended for a single machine and a single analyst working with
a table of up to roughly one hundred thousand rows. The following lists record
precisely what is included and what is not.

\subsection{In scope}

\begin{itemize}
  \item Import of a single-sheet CSV or XLSX file, including day-first dates,
    through an SQL type-inference step.
  \item Discovery of exact and approximate functional dependencies, and of
    single and composite candidate keys, performed by SQL and PL/pgSQL
    functions on the staged data.
  \item Near-duplicate detection with \code{pg\_trgm} and a GIN index, and
    value-level inconsistency fixes recorded in an audit log.
  \item Functional-dependency algebra (attribute closure, minimal cover) and
    Bernstein 3NF synthesis, with generated names and DDL.
  \item A relational metadata schema of 22 tables with constraints, generated
    columns, triggers and reporting views.
  \item Local accounts with salted password hashing and server-side sessions,
    and a PostgreSQL role and pair of schemas for each workspace.
  \item Bundled sample datasets with planted defects: a coaching-institute fee
    register of 3,000 rows, a kirana shop sales register and a hostel mess
    register.
  \item A web interface in a consistent dark and light design system.
  \item Application of a design in one transaction followed by a lossless-join
    check, and an interactive entity--relationship canvas.
  \item A read-only SQL editor, saved queries and run history; a performance
    laboratory; a step-through visualizer; and a natural-language assistant.
\end{itemize}

\subsection{Out of scope}

\begin{itemize}
  \item Database engines other than PostgreSQL.
  \item Multi-sheet workbooks, and files beyond about 20~MB or 100,000 rows.
  \item Write statements from the query editor. The editor is read-only;
    only the application's own flows (apply a design, merge duplicates,
    create an index, scale up data) change user tables.
  \item Step-through of statements that are not a plain \code{SELECT}, or of
    queries dominated by common table expressions or window functions.
  \item Live multi-user collaboration, mobile applications and payments.
  \item Hardening for public hosting on the open internet. The sandbox is
    designed to be sound for a local demonstration, and exposing it publicly
    would be a separate piece of work.
  \item Normal forms beyond 3NF, such as BCNF and 4NF, and the discovery of
    multivalued dependencies.
\end{itemize}

\subsection{Status of the deliverable}

The report describes the system as it stood at the time of writing, and the
status of each capability is marked in the same way throughout. \stdone{} means
implemented and covered by automated tests (212 backend tests and the
frontend test suite pass), \stwip{} means partly implemented, and \stplan{}
means designed but not yet built. The core analytic pipeline from import
through 3NF proposal is built and tested. Applying a design transactionally
with the lossless check, the design canvas and the query sandbox interface are
in progress, and the performance laboratory, the visualizer and the
language-model assistant are planned. No feature is described in this report
as working unless it is marked \stdone.

%==============================================================
\section{Product Scenarios}
\label{sec:scenarios}

Four short scenarios show how the system is expected to be used, who uses it,
and how much of each scenario QLens supports today. A summary of the product
as the user first meets it is shown in \cref{fig:dashboard}.

\screenshot{dashboard-dark}{QLens workspace dashboard (dark theme)}{fig:dashboard}

\subsection{Scenario 1: the coaching-institute owner}

Mrs.\ Verma runs a coaching institute and records all fees in one sheet. She
creates an account, opens the bundled fee register (or uploads her own file),
and presses \emph{Analyze}. QLens infers column types and finds the
dependencies of \cref{sec:motivation}: it reports, for example, that a student
determines a batch, a batch determines a teacher, and a teacher determines a
phone number. Because eight rows spell a student's name differently and one
teacher's earlier number remains in five rows, a few of those dependencies hold
for all rows but a handful; QLens lists the violating rows as inconsistencies
and proposes fixes. After she accepts them, the system proposes four tables:
\tbl{students}, \tbl{batches}, \tbl{teachers} and \tbl{fee\_register}, with
foreign keys between them and a uniqueness constraint on the student and month
of each payment. A teacher's phone number is now stored once. Creating these
tables in her workspace and verifying the migration is the apply step, which
is currently \stwip.

\begin{center}
\small
\begin{tabularx}{\linewidth}{lXl}
\toprule
Step & What QLens does & Status \\
\midrule
Import and type inference & Stages the file, infers SQL types & \stdone \\
Dependency and key discovery & Runs in SQL, about 0.3~s for 3,000 rows & \stdone \\
Duplicate detection and fixes & \code{pg\_trgm} similarity, audit log & \stdone \\
3NF proposal & Bernstein synthesis, names and DDL & \stdone \\
Apply and lossless check & One transaction, \code{EXCEPT ALL} comparison & \stwip \\
\bottomrule
\end{tabularx}
\end{center}

\subsection{Scenario 2: the kirana shop}

A neighbourhood shop keeps a daily sales sheet with the date, item, item
category, supplier, supplier phone, quantity and price. The owner wants to
know which supplier's items sell best, but the same supplier's phone number is
repeated on hundreds of lines and a few lines carry a typing variant of an item
name. Using the bundled kirana dataset, QLens discovers that an item
determines its category and supplier, and that a supplier determines the
phone, and proposes the corresponding split. Once the data is in separate
tables, a question such as total quantity sold per supplier becomes a short
\code{JOIN} with \code{GROUP BY}. The analysis part of this scenario is
\stdone{}; the read-only query editor in which the owner would run that
question is \stwip, and asking it in plain English through the assistant is
\stplan.

\subsection{Scenario 3: the developer migrating legacy CSV}

A developer inherits a decade-old CSV export from a system with no schema
documentation and must load it into an application database. Rather than
guess keys by eye, the developer uploads the file and reads off the candidate
keys and dependencies that actually hold in the data, including composite keys
of up to three columns. The proposed design arrives with generated \code{CREATE
TABLE} statements carrying primary keys, foreign keys and \code{UNIQUE}
constraints, which can be copied into the application's own migrations. The
percentages of violating rows reported next to each approximate dependency tell
the developer where the legacy data cannot be trusted. Discovery and DDL
generation are \stdone{}; the in-database migration and its \code{EXCEPT ALL}
check are \stwip, so a developer today should treat the generated DDL as a
reviewed starting point and not as an automatic migration.

\subsection{Scenario 4: the student learning normalization}

A second-year student is asked to normalize a relation to 3NF and cannot tell
whether the answer is right. Using one of the sample files, the student sees
the dependencies, the candidate keys and the minimal cover derived from them,
and can compare the proposal with a hand-worked solution. Because the
dependencies come from a data set that the student already understands, each
table in the proposal can be traced back to a determinant, which makes the
reasoning about partial and transitive dependencies concrete. Later, the
planned visualizer will show how a \code{SELECT} over the resulting tables
passes through \code{FROM}, \code{WHERE}, \code{GROUP BY}, \code{HAVING} and
\code{ORDER BY}, with row counts at each stage, and the planned performance
laboratory will show the corresponding execution plan. The analysis and the
3NF proposal are \stdone{}; the visualizer and the laboratory are \stplan.

%==============================================================
\section{Organisation of the Report}
\label{sec:organisation}

The remaining chapters follow the order in which the system was specified and
built. The requirements chapter lists the users, the functional and
non-functional requirements and the use cases. The chapter on software
engineering methodology describes how the work was planned and tested. The
system design chapters then present the goals, the architecture, the
relational schema with its entity--relationship diagrams, and the detailed
design of the analysis algorithms, the sandbox, and the user interface, with
screenshots of the features that exist. The report closes with a conclusion
that restates the results honestly against the objectives above, followed by
references and an appendix containing selected SQL listings and test
information.
